\begin{pastexam}[label={px:ntu-108-ee-dsb-13}]{臺大電機所 108 學年度（B）第 13 題}
（複選；每題 10 分，每個選項單獨計分，答對 2 分，答錯倒扣 2 分）

A graph is a pair $(V,E)$, where $V$ is a set of nodes, called vertices, and $E$ is a collection of
pairs of vertices, called edges. There are several ways to represent a graph, among them Adjacency
Matrix and Adjacency List.
\begin{itemize}
\item Adjacency Matrix: Let $A$ be a matrix of size $|V|\times|V|$. $A[u,v]=1$ if there is an edge
between vertex $u$ and $v$, and 0, otherwise.
\item Adjacency List: All vertices connected to a vertex $v$ are listed on an adjacency list for
vertex $v$.
\end{itemize}
Which of the following are correct?
\begin{enumerate}[label=(\Alph*).]
\item Adjacency List is more suitable for dense graph.
\item Using Adjacency Matrix, checking edge between vertices $u$ and $v$ takes constant time.
\item The big advantage of BFS is that it has much lower memory requirement than DFS.
\item BFS is better in finding shortest paths than DFS.
\item Using Adjacency List, the time complexity of topological sort on a directed graph is
$O(|E|+|V|)$.
\end{enumerate}
\end{pastexam}

\begin{pastsoln}
答案是 (B)(D)(E)。

\textbf{(A) 錯}。 相鄰串列的空間是 $\Theta(|V| + |E|)$，相鄰矩陣是 $\Theta(|V|^2)$（資結書
第~21.2~節）。 稀疏圖（$|E|$ 遠小於 $|V|^2$）用串列省空間；稠密圖 $|E| = \Theta(|V|^2)$，兩者
空間同階，矩陣每條邊只佔一個位元、查邊 $O(1)$，又沒有指標的額外負擔，比較適合。 所以串列適合的
是稀疏圖。

\textbf{(B) 對}。 直接讀 $A[u, v]$，一次陣列存取，$O(1)$。 用相鄰串列則要掃 $u$ 的串列，
$O(\deg(u))$。

\textbf{(C) 錯}。 BFS 的佇列裡可能同時放著一整層的頂點，DFS 的堆疊（或遞迴）只放目前這條路徑。
例如星形圖（中心連到其他 $|V| - 1$ 個頂點）從中心出發：BFS 的佇列一下子放進 $|V| - 1$ 個頂點，
DFS 的堆疊深度只有 $2$。 反過來在一條長鏈上 DFS 的堆疊比較深。 最壞情況兩者都是 $\Theta(|V|)$；
硬要比的話，分支多的圖（例如搜尋樹）通常是 DFS 比較省記憶體。 「BFS 的一大優點是記憶體用量
比 DFS 低很多」不成立。

\textbf{(D) 對}。 在無權重的圖（或每條邊等長）上，BFS 依距離由近到遠一層一層發現頂點，得到的
$d[v]$ 就是起點到 $v$ 的最短距離（資結書第~21.3.4~節、定理~21.4）。 DFS 沿一條路一直走下去，
第一次走到 $v$ 的路徑不一定最短，例如三角形 $s$–$a$–$v$ 加上邊 $s$–$v$，DFS 可能先走
$s \to a \to v$，長度 $2$，但最短是 $1$。

\textbf{(E) 對}。 Kahn 演算法先算所有入分支度（掃過每個串列，$O(|V| + |E|)$），之後每個頂點
進出佇列一次、每條邊在它的起點出佇列時被處理一次，總共 $O(|V| + |E|)$（資結書第~22.1~節）。
用 DFS 完成時間排序也是 $O(|V| + |E|)$（資結書第~22.2~節）。
\end{pastsoln}
